% ============================================================
% BookForge — Generated Book Cover
% Compile: xelatex cover.tex (or pdflatex as fallback)
% ============================================================
\documentclass{article}
\usepackage[paperwidth=210mm, paperheight=297mm, margin=0pt]{geometry}
\usepackage{tikz}
\usetikzlibrary{calc,positioning}
\usepackage{eso-pic}

% ── Font setup (xelatex preferred, pdflatex fallback) ──
\usepackage{ifxetex}
\ifxetex
  \usepackage{fontspec}
  \setmainfont{Liberation Sans}[
    BoldFont={Liberation Sans Bold},
    ItalicFont={Liberation Sans Italic}
  ]
  \setsansfont{Liberation Sans}[
    BoldFont={Liberation Sans Bold},
    ItalicFont={Liberation Sans Italic}
  ]
\else
  \usepackage[utf8]{inputenc}
  \usepackage[T1]{fontenc}
  \usepackage{lmodern}
  \usepackage{helvet}
  \renewcommand{\familydefault}{\sfdefault}
\fi

\usepackage{xcolor}
\usepackage{textcomp}
\pagestyle{empty}
\parindent=0pt
\topskip=0pt

% ── Cover colors ──
\definecolor{covbgdeep}{HTML}{011710}
\definecolor{covbgmid}{HTML}{01261A}
\definecolor{covprimary}{HTML}{059669}
\definecolor{covprimarydark}{HTML}{04694A}
\definecolor{covprimarylight}{HTML}{E1F2ED}
\definecolor{covaccent}{HTML}{FFFFFF}
\definecolor{covaccentalt}{HTML}{16A34A}
\definecolor{covtextbright}{HTML}{FFFFFF}
\definecolor{covtextmuted}{HTML}{689386}

% ── Page background (fills any gaps around TikZ picture) ──
\pagecolor{covbgdeep}

\begin{document}%
% [remember picture, overlay] anchors the artwork at the exact page corner
% (needs the 2nd compile pass — compileCover always runs two). This fills the
% page edge-to-edge without the old +5mm paper hack, and unlike an eso-pic
% shipout hook it keeps TikZ transparency working (opacity dies in shipout BG).
\thispagestyle{empty}%
\begin{tikzpicture}[remember picture, overlay,
  shift={(current page.south west)},
  x=1.00000mm, y=1.00000mm]
\clip (0,0) rectangle (210,297);


% === Background ===
\fill[covbgdeep] (0,0) rectangle (210,297);

  % ── Subtle grid ──
  \begin{scope}[opacity=0.035]
    \foreach \x in {0,8,...,210} {
      \draw[covprimary, line width=0.2pt] (\x, 0) -- (\x, 297);
    }
    \foreach \y in {0,8,...,297} {
      \draw[covprimary, line width=0.2pt] (0, \y) -- (210, \y);
    }
  \end{scope}

% === Geometric circles — top-right ===

  % ── Geometric circles ──
  \draw[covprimary, line width=1pt, opacity=0.12]
    (175, 255) circle (75);
  \draw[covprimary, line width=0.5pt, opacity=0.08]
    (175, 255) circle (95);
  \draw[covaccentalt, line width=0.3pt, opacity=0.10]
    (175, 255) circle (110);

% === Circles — bottom-left ===
\draw[covprimary, line width=0.6pt, opacity=0.08] (30, 35) circle (55);
\draw[covprimary, line width=0.3pt, opacity=0.05] (30, 35) circle (70);

% === Glow effects ===
  \fill[covaccentalt, opacity=0.05] (178, 260) circle (45);
  \fill[covprimary, opacity=0.04] (35, 40) circle (35);

% === Background watermark ===
\node[anchor=east, opacity=0.055] at (205, 170) {%
  \fontsize{170}{170}\selectfont\bfseries\sffamily\color{covprimary}%
  60%
};

% === Featured badge circle ===

  % ── Featured badge ──
  \draw[covprimary!50, line width=2.5pt, opacity=0.5]
    (105, 230) circle (26);
  \draw[covaccent, line width=2.8pt, line cap=round]
    (105, 230) ++(90:26) arc (90:430:26);
  % Grot strzałki
  \fill[covaccent]
    ([shift={(80:26)}]105,230) -- ++(70:3.5) -- ++(160:3.5) -- cycle;
  % Tekst wewnątrz
  \node[anchor=center] at (105, 232) {%
    \fontsize{22}{22}\selectfont\bfseries\sffamily\color{covtextbright}%
    60%
  };
  \node[anchor=center] at (105, 221) {%
    \fontsize{6.5}{6.5}\selectfont\sffamily\color{covtextmuted}%
    STRON · KOMPLETNY PRZEWODNIK%
  };

% === Accent line ===
  \fill[covaccent] (15, 200) rectangle (55, 200.8);

% === TITLE ===

\node[anchor=south west] at (15, 179) {%
  \fontsize{28}{31}\selectfont\bfseries\sffamily\color{covtextbright}%
  \mbox{Jak} \mbox{napisać} \mbox{pracę} \mbox{magisterską}%
};
\node[anchor=south west] at (15, 166) {%
  \fontsize{28}{31}\selectfont\bfseries\sffamily\color{covprimary}%
  \mbox{z} \mbox{psychologii}%
};

% === Subtitle ===

  \fill[covaccent] (15, 152) rectangle (70, 152.7);
\node[anchor=north west, text width=170mm] at (15, 149) {%
  \fontsize{16}{20}\selectfont\sffamily\color{covtextmuted}%
  Magisterka z psychologii od A do Z - kompletny przewodnik%
};

% === Topic pillar boxes ===

  % ── Topic pillar boxes ──
  \fill[covprimary, opacity=0.12, rounded corners=2pt]
    (15, 104) rectangle (63, 125);
  \draw[covprimary, line width=0.6pt, rounded corners=2pt]
    (15, 104) rectangle (63, 125);
  \node[anchor=north west] at (18, 123) {%
    \fontsize{8}{8}\selectfont\bfseries\sffamily\color{covaccentalt}%
    OD POMYSŁU DO TEMATU%
  };
  \node[anchor=north west, text width=42mm] at (18, 115) {%
    \fontsize{7}{9}\selectfont\sffamily\color{covtextmuted}%
    Student dowie się, jak przekształcić ogólne zainteresowanie…%
  };
  \fill[covprimary, opacity=0.12, rounded corners=2pt]
    (68, 104) rectangle (116, 125);
  \draw[covprimary, line width=0.6pt, rounded corners=2pt]
    (68, 104) rectangle (116, 125);
  \node[anchor=north west] at (71, 123) {%
    \fontsize{8}{8}\selectfont\bfseries\sffamily\color{covaccentalt}%
    PISANIE ROZDZIAŁÓW%
  };
  \node[anchor=north west, text width=42mm] at (71, 115) {%
    \fontsize{7}{9}\selectfont\sffamily\color{covtextmuted}%
    Czytelnik nauczy się pisać każdy rozdział pracy…%
  };
  \fill[covprimary, opacity=0.12, rounded corners=2pt]
    (121, 104) rectangle (169, 125);
  \draw[covprimary, line width=0.6pt, rounded corners=2pt]
    (121, 104) rectangle (169, 125);
  \node[anchor=north west] at (124, 123) {%
    \fontsize{8}{8}\selectfont\bfseries\sffamily\color{covaccentalt}%
    OBRONA I FINISZ%
  };
  \node[anchor=north west, text width=42mm] at (124, 115) {%
    \fontsize{7}{9}\selectfont\sffamily\color{covtextmuted}%
    Praktyczny przewodnik przez ostatni etap: techniczne wymogi…%
  };

% === Decorative dots (left edge) ===
  \foreach \y in {72,76,...,96} {
    \fill[covprimary, opacity=0.2] (8, \y) circle (0.5);
  }

% === Separator ===
\fill[covprimary, opacity=0.25] (15, 62) rectangle (195, 62.3);

% === Author ===



% === Year ===
\node[anchor=north east] at (195, 55) {%
  \fontsize{13}{13}\selectfont\sffamily\color{covprimary}%
  2026%
};

% === Bottom accent bars ===
  \fill[covaccent] (15, 22) rectangle (40, 22.7);
\fill[covprimary] (43, 22) rectangle (61, 22.7);

% === Bottom-right decoration ===
\fill[covprimary, opacity=0.06] (182, 28) circle (14);
\fill[covaccentalt, opacity=0.04] (188, 22) circle (7);


\end{tikzpicture}%
\end{document}
